\documentclass[11pt]{article}

% -------------------- Packages --------------------
\usepackage[margin=1in]{geometry}
\usepackage{amsmath,amssymb,amsthm,mathtools}
\usepackage{enumitem}
\usepackage{hyperref}
\usepackage{microtype}

% -------------------- Theorem environments --------------------
\theoremstyle{plain}
\newtheorem{theorem}{Theorem}[section]
\newtheorem{lemma}[theorem]{Lemma}
\newtheorem{proposition}[theorem]{Proposition}
\newtheorem{corollary}[theorem]{Corollary}

\theoremstyle{definition}
\newtheorem{definition}[theorem]{Definition}
\newtheorem{assumption}[theorem]{Assumption}
\newtheorem{example}[theorem]{Example}

\theoremstyle{remark}
\newtheorem{remark}[theorem]{Remark}

% -------------------- Notation macros (placeholders) --------------------
\newcommand{\R}{\mathbb{R}}
\newcommand{\E}{\mathbb{E}}
\newcommand{\ri}{\mathrm{ri}}
\newcommand{\cl}{\mathrm{cl}}
\newcommand{\conv}{\mathrm{conv}}
\newcommand{\argmin}{\operatorname*{arg\,min}}
\newcommand{\argmax}{\operatorname*{arg\,max}}
\newcommand{\Row}{\operatorname{Row}}
\newcommand{\col}{\operatorname{col}}
\newcommand{\diag}{\operatorname{diag}}

% -------------------- Title --------------------
\title{Decision-Optimal Single-Scenario Forecasts\\for Two-Stage Stochastic Linear Programs}
\author{}
\date{}

\begin{document}
\maketitle

\tableofcontents

% ============================================================
\section{Setting, notation, and standing assumptions}\label{sec:setting}

\subsection{Two-stage stochastic linear program}\label{subsec:2stage}

Let $(\Omega,\mathcal F,\mathbb P)$ be a probability space and let $\xi$ denote the random second-stage
data.  We work with a two-stage stochastic linear program in (primal) standard form.
The first-stage decision is $z\in\R^{n_1}$ and the second-stage (recourse) decision is $y\in\R^{n_2}$.
The first-stage feasible set is
\begin{equation}\label{eq:Z-def}
  Z := \{\, z\in\R^{n_1}_+ : Az=b \,\},
\end{equation}
where $A\in\R^{m_1\times n_1}$ and $b\in\R^{m_1}$ are deterministic.

A realization $\xi=(q,h,W,T)$ consists of vectors/matrices
\[
q\in\R^{n_2},\qquad h\in\R^{m_2},\qquad W\in\R^{m_2\times n_2},\qquad T\in\R^{m_2\times n_1},
\]
and induces the second-stage linear program
\begin{equation}\label{eq:recourse-LP}
  Q(z,\xi)\;:=\;\min_{y\in\R^{n_2}_+}\ \{\, q^\top y : Wy = h - Tz \,\}.
\end{equation}
The two-stage stochastic program is then
\begin{equation}\label{eq:2stage-SP}
  \min_{z\in Z}\ \Big\{\, c^\top z + \E\big[\, Q(z,\xi)\,\big] \,\Big\},
\end{equation}
where $c\in\R^{n_1}$ is deterministic.  We emphasize that, depending on the application, $\xi$ may contain
any subset of the components $(q,h,W,T)$; the remaining components can be deterministic.

\paragraph{Standard-form convention.}
Writing $Z$ and the recourse constraints in equality$+$nonnegativity form is purely notational: any two-stage
stochastic linear program with polyhedral constraints can be rewritten in this form by introducing slack
variables and splitting free variables into positive and negative parts.  Our assumptions below should be
understood as assumptions on this standard-form representation.

\subsection{Finite support and deterministic equivalent (extensive form)}\label{subsec:extensive}

When $\xi$ has finite support $\{\xi_1,\dots,\xi_N\}$ with probabilities $p_i:=\mathbb P(\xi=\xi_i)>0$,
\[
\xi_i=(q_i,h_i,W_i,T_i), \qquad i\in[N]:=\{1,\dots,N\}, \qquad \sum_{i=1}^N p_i=1,
\]
the expectation in \eqref{eq:2stage-SP} becomes a finite sum and \eqref{eq:2stage-SP} is equivalent to the
deterministic extensive-form linear program
\begin{equation}\label{eq:extensive-form}
\begin{aligned}
  \min_{z,\,y_1,\dots,y_N}\quad & c^\top z + \sum_{i=1}^N p_i\, q_i^\top y_i \\
  \text{s.t.}\quad & Az=b,\quad z\ge 0,\\
  & W_i y_i = h_i - T_i z,\quad y_i\ge 0,\qquad i\in[N].
\end{aligned}
\end{equation}
It will be convenient to write \eqref{eq:extensive-form} in the compact form
\begin{equation}\label{eq:extensive-compact}
  \min_{x\ge 0}\ \bar c^\top x \quad\text{s.t.}\quad \bar A x = \bar b,
\end{equation}
where $x:=(z,y_1,\dots,y_N)$, $\bar c:=(c,p_1q_1,\dots,p_N q_N)$, $\bar b:=(b,h_1,\dots,h_N)$, and
\begin{equation}\label{eq:Abar}
  \bar A :=
  \begin{bmatrix}
    A & 0 & 0 & \cdots & 0 \\
    T_1 & W_1 & 0 & \cdots & 0 \\
    T_2 & 0 & W_2 & \cdots & 0 \\
    \vdots & \vdots & \vdots & \ddots & \vdots \\
    T_N & 0 & 0 & \cdots & W_N
  \end{bmatrix}.
\end{equation}

\subsection{Single-scenario problems and canonical solution maps}\label{subsec:solution-maps}

Fix a \emph{single} scenario $\xi=(q,h,W,T)$ (not necessarily one of the $\xi_i$ in the support).
The corresponding single-scenario linear program is
\begin{equation}\label{eq:single-scenario-LP}
\begin{aligned}
  \min_{z,y}\quad & c^\top z + q^\top y \\
  \text{s.t.}\quad & Az=b,\quad z\ge 0,\\
  & Wy = h - Tz,\quad y\ge 0.
\end{aligned}
\end{equation}
For $\mu>0$, we also consider the log-barrier regularization of \eqref{eq:single-scenario-LP}:
\begin{equation}\label{eq:single-scenario-barrier}
\begin{aligned}
  \min_{z,y}\quad & c^\top z + q^\top y\;-\;\mu\sum_{j=1}^{n_1}\log(z_j)\;-\;\mu\sum_{k=1}^{n_2}\log(y_k)\\
  \text{s.t.}\quad & Az=b,\\
  & Wy = h - Tz,\\
  & z>0,\ y>0.
\end{aligned}
\end{equation}

We also introduce the (multi-scenario) first-stage value functions associated with \eqref{eq:2stage-SP}.
For a given $z\in Z$, define
\begin{equation}\label{eq:G-def}
  G(z)\;:=\;c^\top z + \sum_{i=1}^N p_i\,Q(z,\xi_i),
\end{equation}
and for $\mu>0$ define the barrier-regularized recourse functions
\begin{equation}\label{eq:Qmu-def}
  Q_\mu(z,\xi_i)\;:=\;\min_{y>0}\ \Big\{\, q_i^\top y - \mu\sum_{k=1}^{n_2}\log(y_k)\ :\ W_i y = h_i - T_i z \,\Big\},
\end{equation}
together with the barrier-regularized first-stage value function
\begin{equation}\label{eq:Gmu-def}
  G_\mu(z)\;:=\;c^\top z - \mu\sum_{j=1}^{n_1}\log(z_j) + \sum_{i=1}^N p_i\,Q_\mu(z,\xi_i),
\end{equation}
defined on $Z\cap\{z>0\}$ whenever the second-stage terms are well-posed.

\begin{definition}[Canonical single-scenario decision map]\label{def:canonical-map}
Fix a scenario $\xi=(q,h,W,T)$ and consider the single-scenario feasible set
\[
  \mathcal P(\xi)\;:=\;\big\{(z,y)\in\R^{n_1}\times\R^{n_2}:\ Az=b,\ Wy=h-Tz,\ z\ge 0,\ y\ge 0\big\}.
\]

\smallskip
\noindent\textbf{Barrier map.}
For $\mu>0$, let $(z_\mu^*(\xi),y_\mu^*(\xi))$ be the (unique) minimizer of the strictly convex barrier problem
\eqref{eq:single-scenario-barrier}.  We call $z_\mu^*(\xi)\in\R^{n_1}$ the \emph{canonical barrier decision}
induced by $\xi$.

\smallskip
\noindent\textbf{Unsmooth canonical map.}
Let
\[
  \mathcal S(\xi)\;:=\;\argmin\big\{\, c^\top z + q^\top y : (z,y)\in\mathcal P(\xi)\,\big\}
\]
be the (possibly set-valued) optimal solution set of \eqref{eq:single-scenario-LP}.  Define the index set of
coordinates that are not identically zero on $\mathcal S(\xi)$ by
\[
  \mathcal J(\xi)\;:=\;\Big\{\ell:\ \exists (z,y)\in\mathcal S(\xi)\ \text{with}\ (z,y)_\ell>0\Big\}.
\]
The \emph{analytic center of the optimal face} is the unique point $(z^*(\xi),y^*(\xi))$ solving
\begin{equation}\label{eq:analytic-center-opt-face}
  \max_{(z,y)\in\mathcal S(\xi)}\ \sum_{\ell\in\mathcal J(\xi)} \log\big((z,y)_\ell\big),
\end{equation}
where the maximization is taken over the relative interior of $\mathcal S(\xi)$ in its affine hull.  We call
$z^*(\xi)$ the \emph{canonical (unsmoothed) decision} induced by $\xi$.
\end{definition}

\subsection{Decision-optimal single-scenario forecasts}\label{subsec:dosf}

We now formalize what it means for a \emph{single} scenario to be ``decision-optimal'' for the stochastic
problem: it should reproduce the \emph{canonical} optimal first-stage decision.

Let $z_\mu^\star$ denote the (unique) first-stage component of the optimal solution of the barrier-smoothed
extensive form (defined later in Section~\ref{sec:barrier}, but already unambiguous under our
assumptions).  Let $z^\star$ denote the first-stage component of the analytic center of the optimal face of the
unsmoothed extensive form \eqref{eq:extensive-form}.

\begin{definition}[Decision-optimal single-scenario forecast]\label{def:dosf}
A scenario $\hat\xi$ (not necessarily in the support $\{\xi_1,\dots,\xi_N\}$) is called
\begin{enumerate}[leftmargin=2.2em]
\item \emph{decision-optimal for the barrier problem at parameter $\mu>0$} if
\[
  z_\mu^*(\hat\xi)\;=\;z_\mu^\star.
\]
\item \emph{decision-optimal for the unsmoothed problem} if
\[
  z^*(\hat\xi)\;=\;z^\star.
\]
\end{enumerate}
\end{definition}

The use of the canonical maps $z_\mu^*(\cdot)$ and $z^*(\cdot)$ is essential: a single scenario can make the
stochastic optimum \emph{one} of many optimal first-stage decisions without actually encoding any meaningful
decision information.

\subsection{A cautionary example: vacuous ``optimal'' scenarios}\label{subsec:vacuous}

\begin{example}[A forecast may be decision-irrelevant]\label{ex:vacuous}
Consider a two-stage problem with a one-dimensional first-stage decision $z\in[0,1]$ written in standard
form via $z+s=1$ with $(z,s)\ge 0$.  There is a single second-stage variable $y\ge 0$ constrained by $y=z$
(i.e., $y-z=0$).  The second-stage cost is random: $\xi=q\in\R_+$ and $Q(z,q)=\min\{qy: y=z,\ y\ge 0\}=qz$.
There is no first-stage cost ($c=0$), hence the stochastic objective is $\E[q]\,z$, and the unique optimal
first-stage decision is $z^\star=0$.

Now consider the single-scenario problem with the \emph{forecast} $\hat q=0$.  Then the single-scenario
objective is identically zero for every feasible $z\in[0,1]$, i.e., every first-stage decision is optimal.
In particular, the forecast $\hat q=0$ \emph{contains} the true optimum $z^\star=0$ in its optimal solution set,
but it is nonetheless decision-irrelevant: it cannot distinguish $z^\star$ from any other feasible decision.
Under our canonical selection rule, $z^*(\hat q)$ is the analytic center of $[0,1]$, namely $z^*(\hat q)=1/2$,
which is \emph{not} equal to $z^\star$.

This illustrates why we insist on Definition~\ref{def:dosf}: decision-optimality requires reproducing the
canonical optimal decision, not merely including it among many optima.
\end{example}

\subsection{Standing assumptions}\label{subsec:assumptions}

We state the standing assumptions that will be used throughout the remainder of the paper.

\begin{assumption}[Finite support]\label{ass:finite-support}
The random element $\xi$ has finite support $\{\xi_1,\dots,\xi_N\}$ with probabilities $p_i>0$ for all $i\in[N]$.
\end{assumption}

\begin{assumption}[No redundant constraints (full row rank)]\label{ass:full-row-rank}
The extensive-form equality constraint matrix $\bar A$ defined in \eqref{eq:Abar} has full row rank.
Equivalently, the extensive form \eqref{eq:extensive-form} has no redundant equality constraints.
\end{assumption}

\begin{assumption}[Bounded and relatively complete recourse]\label{ass:recourse}
For every first-stage decision $z\in Z$ and every scenario $i\in[N]$, the recourse problem
\[
  Q(z,\xi_i)=\min\{\, q_i^\top y : W_i y = h_i - T_i z,\ y\ge 0 \,\}
\]
is feasible and has a finite optimal value (in particular, it is not unbounded below).
\end{assumption}

\begin{assumption}[Existence of a solution]\label{ass:existence}
The extensive-form problem \eqref{eq:extensive-form} has at least one optimal solution (i.e., its optimal
value is finite and attained).
\end{assumption}

\begin{assumption}[Strict feasibility / nonempty relative interior]\label{ass:ri-nonempty}
The feasible region of the extensive form has nonempty relative interior:
there exist $(\tilde z,\tilde y_1,\dots,\tilde y_N)$ satisfying
\[
  A\tilde z=b,\qquad W_i\tilde y_i = h_i - T_i\tilde z\ \ (i\in[N]),\qquad \tilde z>0,\ \tilde y_i>0\ \ (i\in[N]).
\]
\end{assumption}

\begin{assumption}[Recession cone / barrier well-posedness]\label{ass:recession}
Let $\mathcal R:=\{d\in\R^{n_1+Nn_2}_+ : \bar A d=0\}$ be the recession cone of the extensive-form feasible set.
Then
\[
  \bar c^\top d \;>\;0 \qquad \text{for all nonzero } d\in\mathcal R,
\]
where $\bar c=(c,p_1q_1,\dots,p_Nq_N)$ as in \eqref{eq:extensive-compact}.
\end{assumption}

\begin{remark}[On restrictiveness and SAA]\label{rem:saa}
Assumption~\ref{ass:finite-support} is primarily a notational convenience: it allows us to work directly with
the extensive form \eqref{eq:extensive-form}.  For general (non-finite) distributions, one can replace the
expectation in \eqref{eq:2stage-SP} with a sample-average approximation (SAA) and apply the results of this
paper to the resulting finitely-supported empirical problem.  In that setting, our existence statements yield
decision-optimal (for the SAA) single-scenario forecasts; out-of-sample guarantees can then be obtained by
combining SAA consistency with standard stability arguments (e.g., $\varepsilon$-optimality under sufficiently
large samples).
\end{remark}

\begin{remark}[On full row rank and interior feasibility]\label{rem:rank-ri}
Assumption~\ref{ass:full-row-rank} can always be enforced by deleting redundant equalities from the
standard-form representation; it is used to avoid degeneracy in the barrier KKT system and to guarantee
uniqueness properties for the barrier solutions.

Assumption~\ref{ass:ri-nonempty} is equivalent to the existence of a \emph{single} strictly feasible first-stage
decision $\tilde z\in\ri(Z)$ such that, for every scenario $i$, the second-stage feasible region
$\{y\ge 0: W_i y = h_i - T_i\tilde z\}$ has nonempty relative interior (i.e., admits a strictly positive $y$).
\end{remark}

\begin{remark}[On standardity of barrier assumptions]\label{rem:barrier-standard}
Assumptions~\ref{ass:ri-nonempty}--\ref{ass:recession} are precisely the well-posedness conditions needed to
treat the barrier-regularized problems as bona fide convex optimization problems: feasibility of the strict
interior and exclusion of recession directions along which the barrier objective can diverge to $-\infty$.
While these assumptions may appear specific, they are standard in interior-point theory and can typically be
ensured by mild reformulations (e.g., adding slack variables, eliminating redundant equalities, and working in
a standard-form representation).
\end{remark}

% ============================================================
% Section 2: Log-barrier smoothing
% ============================================================

\section{Log-barrier smoothing}\label{sec:barrier}

This section records the log-barrier facts used later to construct \emph{decision-optimal single-scenario
forecasts}. We work throughout under the standing assumptions stated in Section~\ref{sec:setting}.

\subsection{Barrier-smoothed recourse and its dual}\label{subsec:barrier:recourse}

Fix a scenario $\xi=(q,W,T,h)$ and a barrier parameter $\mu>0$. Recall the (unsmoothed) second-stage
recourse value for a first-stage decision $z$ (see \eqref{eq:recourse-LP}):
\[
Q(z,\xi)\;:=\;\min_{y\in\mathbb{R}^{n_2}_+}\Bigl\{q^\top y:\;Wy=h-Tz\Bigr\},
\]
with the convention $Q(z,\xi)=+\infty$ if the equality system has no feasible $y\ge 0$.

\begin{definition}[Log-barrier smoothed recourse]\label{def:Qmu}
For $\mu>0$ we define the log-barrier smoothed recourse
\begin{equation}\label{eq:recourse-barrier}
Q_\mu(z,\xi)\;:=\;\min_{y\in\mathbb{R}^{n_2}_{++}}\Bigl\{q^\top y-\mu\sum_{k=1}^{n_2}\log(y_k):\;Wy=h-Tz\Bigr\},
\end{equation}
again taking $Q_\mu(z,\xi)=+\infty$ if the equality system has no strictly positive solution.
For a vector $v\in\mathbb{R}^{n_2}_{++}$ we use $v^{-1}$ to denote the componentwise reciprocal.
\end{definition}

\begin{proposition}[Dual of the barrier recourse]\label{prop:barrier-dual}
Fix $\xi=(q,W,T,h)$ and $\mu>0$, and write $r(z):=h-Tz$. Assume \eqref{eq:recourse-barrier} is feasible
for the given $z$ (so $Q_\mu(z,\xi)<+\infty$). Then strong duality holds and
\begin{equation}\label{eq:barrier-dual-recourse}
Q_\mu(z,\xi)
=\mu n_2\bigl(1-\log\mu\bigr)
+\max_{\lambda\in\mathbb{R}^m:\;W^\top\lambda<q}
\Bigl\{r(z)^\top\lambda+\mu\sum_{k=1}^{n_2}\log\bigl(q-W^\top\lambda\bigr)_k\Bigr\}.
\end{equation}
Moreover, the primal minimizer in \eqref{eq:recourse-barrier} is unique; we denote it by
$y_\mu^\star(z;\xi)$. Define the (possibly set-valued) dual maximizer set
\[
\Lambda_\mu(z;\xi)
:= \argmax_{\lambda\in\mathbb{R}^m:\,W^\top\lambda<q}
\left\{ r(z)^\top\lambda + \mu\sum_{k=1}^{n_2}\log\bigl((q-W^\top\lambda)_k\bigr)\right\}.
\]
Then $\Lambda_\mu(z;\xi)\neq\emptyset$ and is convex. If $W$ has full row rank, then the maximizer
in \eqref{eq:barrier-dual-recourse} is unique and we write
$\Lambda_\mu(z;\xi)=\{\lambda_\mu^\star(z;\xi)\}$.
\end{proposition}

\begin{proof}
Consider the Lagrangian for \eqref{eq:recourse-barrier} with multiplier $\lambda\in\mathbb{R}^m$ for
the equality $Wy=r(z)$:
\[
\mathcal{L}(y,\lambda)
= q^\top y-\mu\sum_{k=1}^{n_2}\log(y_k)+\lambda^\top\bigl(r(z)-Wy\bigr)
= r(z)^\top\lambda + \sum_{k=1}^{n_2}\Bigl((q-W^\top\lambda)_k\,y_k-\mu\log y_k\Bigr).
\]
For fixed $\lambda$, the infimum over $y\in\mathbb{R}^{n_2}_{++}$ is finite if and only if
$q-W^\top\lambda\in\mathbb{R}^{n_2}_{++}$; in that case the minimizer is componentwise
\[
y_k=\frac{\mu}{(q-W^\top\lambda)_k},\qquad k=1,\dots,n_2.
\]
Substituting this minimizer yields
\[
\begin{aligned}
\inf_{y>0}\mathcal{L}(y,\lambda)
&= r(z)^\top\lambda + \sum_{k=1}^{n_2}\Bigl(\mu-\mu\log\mu+\mu\log\bigl(q-W^\top\lambda\bigr)_k\Bigr)\\
&= r(z)^\top\lambda+\mu n_2(1-\log\mu)
+\mu\sum_{k=1}^{n_2}\log\bigl(q-W^\top\lambda\bigr)_k.
\end{aligned}
\]
Maximizing over $\lambda$ with $W^\top\lambda<q$ gives \eqref{eq:barrier-dual-recourse}.
Feasibility of \eqref{eq:recourse-barrier} implies the usual constraint qualification for equality constraints
on an open domain, so strong duality holds and the maximum is attained.

Finally, the primal objective in \eqref{eq:recourse-barrier} is strictly convex on $\mathbb{R}^{n_2}_{++}$,
hence the primal minimizer $y_\mu^\star(z;\xi)$ is unique. The dual objective in \eqref{eq:barrier-dual-recourse}
is the sum of an affine function and a strictly concave function of $W^\top\lambda$; if $W$ has full row rank,
the map $\lambda\mapsto W^\top\lambda$ is injective, hence the dual objective is strictly concave in $\lambda$,
and the maximizer $\lambda_\mu^\star(z;\xi)$ is unique.
Even when $W$ is not full row rank, strict concavity in $W^\top\lambda$ implies that $W^\top\lambda$ is the same
for all $\lambda\in\Lambda_\mu(z;\xi)$ (equivalently, $q-W^\top\lambda=\mu(y_\mu^\star(z;\xi))^{-1}$ is unique).
\end{proof}

\begin{remark}[Barrier duality as ``barrier of the dual'']\label{rem:barrier-of-dual}
For the unsmoothed recourse \eqref{eq:recourse-LP}, the dual is
\[
Q(z,\xi)=\max_{\lambda\in\mathbb{R}^m}\Bigl\{r(z)^\top\lambda:\;W^\top\lambda\le q\Bigr\}.
\]
Comparing with \eqref{eq:barrier-dual-recourse}, the log-barrier smoothed primal recourse is (up to the
constant $\mu n_2(1-\log\mu)$) the log-barrier regularization of the dual constraints $W^\top\lambda\le q$,
via the term $\mu\sum_k\log\bigl(q-W^\top\lambda\bigr)_k$.
\end{remark}

\subsection{KKT conditions and continuous dependence}\label{subsec:barrier:kkt}

\begin{proposition}[Second-stage KKT system]\label{prop:second-stage-kkt}
Assume $Q_\mu(z,\xi)<+\infty$. Let $y_\mu^\star(z;\xi)$ be the unique primal minimizer of
\eqref{eq:recourse-barrier} and let $\Lambda_\mu(z;\xi)$ be the (nonempty) dual maximizer set defined in
Proposition~\ref{prop:barrier-dual}. Then:
\begin{enumerate}
\item For any $\lambda\in\Lambda_\mu(z;\xi)$, the pair $(y_\mu^\star(z;\xi),\lambda)$ satisfies the KKT system
\begin{equation}\label{eq:KKT-second-stage}
\begin{aligned}
Wy &= h-Tz,\\
q-W^\top\lambda &= \mu\,y^{-1},
\end{aligned}
\qquad\text{with}\quad y\in\mathbb{R}^{n_2}_{++},\;q-W^\top\lambda\in\mathbb{R}^{n_2}_{++}.
\end{equation}
\item Conversely, any pair $(y,\lambda)$ satisfying \eqref{eq:KKT-second-stage} is primal--dual optimal
for \eqref{eq:recourse-barrier}--\eqref{eq:barrier-dual-recourse}.
\end{enumerate}
In particular, $q-W^\top\lambda=\mu(y_\mu^\star(z;\xi))^{-1}$ is unique, so $W^\top\lambda$ is the same
for all $\lambda\in\Lambda_\mu(z;\xi)$.
\end{proposition}

\begin{proof}
Because the objective in \eqref{eq:recourse-barrier} is strictly convex on $\mathbb{R}^{n_2}_{++}$, the primal
minimizer $y_\mu^\star(z;\xi)$ is unique. Strong duality and attainment hold as in
Proposition~\ref{prop:barrier-dual}. Stationarity of the Lagrangian with respect to $y$ yields
$q-W^\top\lambda-\mu y^{-1}=0$, and primal feasibility gives $Wy=h-Tz$, i.e.\ \eqref{eq:KKT-second-stage}.
Conversely, any $(y,\lambda)$ satisfying \eqref{eq:KKT-second-stage} is primal feasible and satisfies
stationarity; since the problem is strictly convex with equality constraints on an open domain, KKT
conditions are sufficient, hence $(y,\lambda)$ is primal--dual optimal.
\end{proof}

\begin{proposition}[Continuity of barrier minimizers for fixed $\mu>0$]\label{prop:IFT}
Fix $\mu>0$. Consider the one-scenario barrier problem \eqref{eq:barrier-one-scenario} with data
$\xi=(q,W,T,h)$. On any set of scenarios for which the barrier problem is feasible and attains a finite
minimum, the minimizer $(z_\mu^\star(\xi),y_\mu^\star(\xi))$ is unique and the mapping
\[
\xi \longmapsto (z_\mu^\star(\xi),y_\mu^\star(\xi))
\]
is continuous.
\end{proposition}

\begin{proof}
For fixed $\mu>0$, the barrier objective is continuous in $(z,y,\xi)$ on its domain and strictly convex
in $(z,y)$ on the affine feasible set. Under the stated feasibility and attainment assumptions, the
argmin is nonempty and compact-valued (minimizers lie in a bounded level set), hence Berge's
maximum theorem gives upper hemicontinuity of the solution map. Uniqueness of the minimizer
upgrades this to continuity.
\end{proof}

\subsection{Subgradients and first-stage optimality}\label{subsec:barrier:grad}

We now turn to the smoothed \emph{first-stage} objective.
For the finite support $\{(\xi_i,p_i)\}_{i=1}^N$ introduced in Section~\ref{sec:setting}, define
\[
G_\mu(z)
:= c^\top z - \mu\sum_{j=1}^{n_1}\log(z_j) + \sum_{i=1}^N p_i\,Q_\mu(z,\xi_i),
\]
with domain consisting of $z$ satisfying the first-stage constraints (e.g. $Az=b$) and for which each
recourse problem $Q_\mu(z,\xi_i)$ is feasible.
When $\mu=0$, $G_0(z)=c^\top z+\sum_i p_iQ(z,\xi_i)$ is the usual risk-neutral first-stage cost.

\begin{proposition}[Subgradients of $Q_\mu$ and of $G_\mu$]\label{prop:grad}
Fix $\mu>0$. Assume $Q_\mu(z,\xi_i)<+\infty$ for the $z$ under consideration. Let
$\Lambda_\mu(z;\xi_i)$ be the dual maximizer set from Proposition~\ref{prop:barrier-dual} for scenario
$\xi_i$. Then $Q_\mu(\cdot,\xi_i)$ is convex and its subdifferential at $z$ satisfies
\begin{equation}\label{eq:subgrad-Qmu}
\partial_z Q_\mu(z,\xi_i)
=
\conv\left\{-T_i^\top\lambda:\ \lambda\in\Lambda_\mu(z;\xi_i)\right\}.
\end{equation}
In particular, for any $\lambda\in\Lambda_\mu(z;\xi_i)$ we have $-T_i^\top\lambda\in\partial_z Q_\mu(z,\xi_i)$.
If $W_i$ has full row rank, then $\Lambda_\mu(z;\xi_i)$ is a singleton and $Q_\mu(\cdot,\xi_i)$ is
differentiable with $\nabla_z Q_\mu(z,\xi_i)=-T_i^\top\lambda_\mu^\star(z;\xi_i)$.

Consequently, $G_\mu$ is convex and
\begin{equation}\label{eq:subgrad-Gmu}
\partial G_\mu(z)
=
c-\mu z^{-1} + \sum_{i=1}^N p_i\,\partial_z Q_\mu(z,\xi_i).
\end{equation}
\end{proposition}

\begin{proof}
Using the dual representation \eqref{eq:barrier-dual-recourse}, the only dependence on $z$ is through
the affine term $(h_i-T_i z)^\top\lambda$. For each fixed $\lambda$, the derivative of that affine term is
$-T_i^\top\lambda$. Danskin's theorem (subdifferential form) yields \eqref{eq:subgrad-Qmu}. Summing with
weights $p_i$ and adding the explicit derivative of $c^\top z-\mu\sum_j\log z_j$ gives \eqref{eq:subgrad-Gmu}.
\end{proof}

\begin{proposition}[KKT conditions for the smoothed first-stage decision]\label{prop:first-stage-KKT}
Fix $\mu>0$ and consider the smoothed stochastic program
\begin{equation}\label{eq:stoch-mu}
\min_{z\in\mathbb{R}^{n_1}_{++}}\;\Bigl\{G_\mu(z):\;Az=b\Bigr\}.
\end{equation}
Assume it attains its minimum (standing assumptions). Let $z_\mu^\star$ denote its unique minimizer.
Then there exist a multiplier $\nu_\mu^\star$ for the equality constraint $Az=b$ and vectors
$\lambda_{\mu}^{\star,i}\in\Lambda_\mu(z_\mu^\star;\xi_i)$, $i=1,\dots,N$, such that
\begin{equation}\label{eq:KKT-first-stage}
Az_\mu^\star=b,\qquad
c-\mu\,(z_\mu^\star)^{-1}+A^\top\nu_\mu^\star-\sum_{i=1}^N p_i\,T_i^\top\,\lambda_{\mu}^{\star,i}=0.
\end{equation}
\end{proposition}

\begin{proof}
By Proposition~\ref{prop:grad}, $G_\mu$ is convex on its domain and its subdifferential is given by
\eqref{eq:subgrad-Gmu}. The feasible set $\{z>0:Az=b\}$ is affine. Standard first-order conditions for
equality-constrained convex minimization imply the existence of $\nu_\mu^\star$ such that
$0\in\partial G_\mu(z_\mu^\star)+A^\top\nu_\mu^\star$ together with feasibility $Az_\mu^\star=b$.
Expanding $\partial G_\mu(z_\mu^\star)$ via \eqref{eq:subgrad-Gmu} and selecting
$\lambda_{\mu}^{\star,i}\in\Lambda_\mu(z_\mu^\star;\xi_i)$ yields \eqref{eq:KKT-first-stage}.
Uniqueness of $z_\mu^\star$ follows from strict convexity of $-\sum_j\log z_j$ on the affine feasible set.
\end{proof}

\subsection[As mu to 0: central-path convergence]{As $\mu\downarrow 0$: central-path convergence}\label{subsec:barrier:convergence}

We next relate the barrier-smoothed solutions to their unsmoothed counterparts. We state the result in the
single-scenario form; the extensive-form version follows by applying the same argument to the extensive-form
LP.

Fix a scenario $\xi=(q,W,T,h)$ and consider the one-scenario LP
\begin{equation}\label{eq:LP-one-scenario}
\min_{z\ge 0,\;y\ge 0}\Bigl\{c^\top z+q^\top y:\;Az=b,\;Wy=h-Tz\Bigr\}.
\end{equation}
Let $\mathcal{F}(\xi)$ denote its feasible set and $\mathcal{S}(\xi)$ its optimal solution set, assumed nonempty.
Under our standing recession-cone assumption, $\mathcal{S}(\xi)$ is a bounded polyhedron.
Let $\mathcal{S}^\circ(\xi):=\mathcal{S}(\xi)\cap\ri(\mathcal{F}(\xi))$ denote the relative interior of the optimal face.

\begin{definition}[Analytic-center selection for LP optima]\label{def:ac}
Assume $\mathcal{S}^\circ(\xi)\neq\emptyset$ (standing assumptions imply this). The \emph{analytic center}
of the optimal face is the unique point $(z^{\mathrm{ac}}(\xi),y^{\mathrm{ac}}(\xi))\in\mathcal{S}^\circ(\xi)$ solving
\begin{equation}\label{eq:ac}
(z^{\mathrm{ac}}(\xi),y^{\mathrm{ac}}(\xi))
\in\argmax_{(z,y)\in\mathcal{S}^\circ(\xi)}
\Bigl\{\sum_{j=1}^{n_1}\log(z_j)+\sum_{k=1}^{n_2}\log(y_k)\Bigr\}.
\end{equation}
We define the \emph{canonical} first-stage decision of the one-scenario LP as
\[
z^\star(\xi):=z^{\mathrm{ac}}(\xi).
\]
\end{definition}

\begin{theorem}[Barrier limit selects the analytic center]\label{thm:central-path-limit}
Fix $\xi=(q,W,T,h)$ and assume the one-scenario LP \eqref{eq:LP-one-scenario} is feasible, has a finite
optimal value, and that $\ri(\mathcal{F}(\xi))\neq\emptyset$ and $\mathcal{S}^\circ(\xi)\neq\emptyset$.
For each $\mu>0$, consider the barrier-regularized one-scenario problem
\begin{equation}\label{eq:barrier-one-scenario}
\min_{z\in\mathbb{R}^{n_1}_{++},\;y\in\mathbb{R}^{n_2}_{++}}
\Bigl\{c^\top z+q^\top y-\mu\sum_{j=1}^{n_1}\log z_j-\mu\sum_{k=1}^{n_2}\log y_k:\;Az=b,\;Wy=h-Tz\Bigr\},
\end{equation}
and let $(z_\mu^\star(\xi),y_\mu^\star(\xi))$ denote its unique minimizer.
Then
\[
(z_\mu^\star(\xi),y_\mu^\star(\xi))\xrightarrow[\mu\downarrow 0]{}\bigl(z^{\mathrm{ac}}(\xi),y^{\mathrm{ac}}(\xi)\bigr),
\qquad\text{in particular}\quad z_\mu^\star(\xi)\to z^\star(\xi).
\]
\end{theorem}

This is a standard central-path result from interior-point theory; see, e.g., Nesterov and Nemirovskii (1994) or Wright (1997).

% ============================================================
% Section 3: Existence of decision-optimal single-scenario forecasts
% ============================================================

\section{Existence of decision-optimal single-scenario forecasts}\label{sec:existence}

In this section we prove that, for several ``noise regimes'', there exist \emph{single-scenario forecasts}
$\hat\xi$ for which solving the one-scenario problem returns exactly the (canonical) optimal first-stage
decision of the original stochastic program. We treat first the $\mu$-smoothed setting (where the canonical
solution is the unique log-barrier minimizer), and then show how to pass to the unsmoothed LP via
$\mu\downarrow 0$.

\subsection[A technical lemma: translating q translates the barrier dual]{A technical lemma: translating $q$ translates the barrier dual}\label{subsec:q-translation}

\begin{lemma}[$q$-translation for the barrier dual]\label{lem:q-translation}
Fix $(W,T,h,z)$ and $\mu>0$. Let $q\in\mathbb{R}^{n_2}$ be such that the dual domain
$\{\lambda:W^\top\lambda<q\}$ is nonempty, and let $\Delta\in\mathbb{R}^m$.
Define $q^+:=q+W^\top\Delta$.
Then the barrier-dual maximizer sets satisfy
\begin{equation}\label{eq:q-translation}
\argmax_{\lambda:\;W^\top\lambda<q^+}
\Bigl\{(h-Tz)^\top\lambda+\mu\sum_{k=1}^{n_2}\log\bigl(q^+-W^\top\lambda\bigr)_k\Bigr\}
=
\Bigl(\argmax_{\lambda:\;W^\top\lambda<q}
\{\cdots\}\Bigr)+\Delta,
\end{equation}
i.e.\ maximizers translate by $\Delta$. In particular, if $W$ has full row rank so that the maximizer is unique,
then
\[
\lambda_\mu^\star(z;q^+,W,T,h)=\lambda_\mu^\star(z;q,W,T,h)+\Delta.
\]
\end{lemma}

\begin{proof}
Consider the barrier dual objective for $q^+$:
\[
\Psi_{q^+}(\lambda)
:=
(h-Tz)^\top\lambda+\mu\sum_{k=1}^{n_2}\log\bigl(q^+-W^\top\lambda\bigr)_k,
\qquad\text{dom}(\Psi_{q^+})=\{\lambda:W^\top\lambda<q^+\}.
\]
With the change of variables $\lambda=\tilde\lambda+\Delta$, we have
$W^\top\lambda<q^+ \iff W^\top\tilde\lambda<q$, and
\[
\Psi_{q^+}(\tilde\lambda+\Delta)
=
(h-Tz)^\top\tilde\lambda+(h-Tz)^\top\Delta
+\mu\sum_{k=1}^{n_2}\log\bigl(q-W^\top\tilde\lambda\bigr)_k.
\]
The term $(h-Tz)^\top\Delta$ is constant in $\tilde\lambda$, so the maximizers in $\tilde\lambda$ are exactly those
for $q$, and the stated translation of maximizers follows.
\end{proof}

\subsection{Decision-optimal forecasts for the barrier-smoothed problem}\label{subsec:existence:mu}

Throughout this subsection, fix $\mu>0$ and let $z_\mu^\star$ be the unique minimizer of the smoothed
stochastic program \eqref{eq:stoch-mu}. For each scenario $i$ and each $z$, let
$\Lambda_\mu^i(z):=\Lambda_\mu(z;\xi_i)$ denote the (possibly set-valued) dual maximizer set in
\eqref{eq:barrier-dual-recourse} for the barrier recourse $Q_\mu(z,\xi_i)$. When $W_i$ is rank-deficient,
$\Lambda_\mu^i(z)$ need not be a singleton.

In the constructions below we only need dual multipliers at the stochastic optimum $z_\mu^\star$.
We fix a canonical choice as follows. Consider the extensive-form barrier problem equivalent to
\eqref{eq:stoch-mu}:
\[
\min_{\substack{z\in\mathbb{R}^{n_1}_{++}\\ y_i\in\mathbb{R}^{n_2}_{++}}}
\ c^\top z - \mu\sum_{j=1}^{n_1}\log z_j
+\sum_{i=1}^N p_i\left(q_i^\top y_i-\mu\sum_{k=1}^{n_2}\log (y_i)_k\right)
\quad\text{s.t.}\quad
Az=b,\ \ W_i y_i = h_i - T_i z\ (i\in[N]).
\]
Let $(z_\mu^\star,y_{\mu,1}^\star,\dots,y_{\mu,N}^\star)$ be its unique minimizer. Under
Assumption~\ref{ass:full-row-rank}, the equality-constraint multipliers are unique; let $\nu_\mu^\star$
be the multiplier for $Az=b$ and let $\widehat\lambda_{\mu}^{\star,i}$ be the multiplier for
$W_i y_i = h_i - T_i z$. Define the scaled multipliers
\[
\lambda_{\mu}^{\star,i} := \widehat\lambda_{\mu}^{\star,i}/p_i \qquad (i\in[N]).
\]
Then for each $i$, $(y_{\mu,i}^\star,\lambda_{\mu}^{\star,i})$ satisfies the second-stage KKT system
\eqref{eq:KKT-second-stage} at $z=z_\mu^\star$, hence $\lambda_{\mu}^{\star,i}\in\Lambda_\mu^i(z_\mu^\star)$.

\subsubsection[Noise in (q,h) with fixed (W,T), and q* chosen a priori]{Noise in $(q,h)$ with fixed $(W,T)$, and \texorpdfstring{$q^\star$}{q*} chosen \emph{a priori}}\label{subsec:qh-fixedWT}

Assume in this subsection that $W_i\equiv W$ and $T_i\equiv T$ are deterministic and fixed across scenarios,
while $(q_i,h_i)$ may vary with $i$.

\begin{theorem}[A decision-optimal $(q^\star,h_\mu^\star)$ under fixed $(W,T)$]\label{thm:qh-fixedWT-barrier}
Assume $W$ and $T$ are fixed across scenarios, and let $z_\mu^\star$ be the unique minimizer of the smoothed
stochastic program \eqref{eq:stoch-mu}. Define the averaged second-stage dual vector
\[
\bar\lambda_\mu:=\sum_{i=1}^N p_i\,\lambda_{\mu}^{\star,i}\in\mathbb{R}^m,
\]
and fix a parameter $\delta>0$ (chosen a priori). Define the componentwise maximum
\[
q^{\max}_k:=\max_{i=1,\dots,N}(q_i)_k,\qquad k=1,\dots,n_2,
\]
and set
\begin{equation}\label{eq:qmax}
(q^\star)_k:= q^{\max}_k + \delta\,|q^{\max}_k|,\qquad k=1,\dots,n_2,
\end{equation}
where $|\cdot|$ denotes absolute value (componentwise).
Then $q^\star-W^\top\bar\lambda_\mu\in\mathbb{R}^{n_2}_{++}$, and the quantities
\[
y_\mu^\star:=\mu\,(q^\star-W^\top\bar\lambda_\mu)^{-1}\in\mathbb{R}^{n_2}_{++},
\qquad
h_\mu^\star:=Tz_\mu^\star+Wy_\mu^\star\in\mathbb{R}^m
\]
define a single-scenario forecast $\hat\xi_\mu^\star:=(q^\star,W,T,h_\mu^\star)$ such that
\begin{equation}\label{eq:decision-optimal-qh}
z_\mu^\star = z_\mu^\star(\hat\xi_\mu^\star),
\end{equation}
i.e.\ the one-scenario barrier problem with scenario $\hat\xi_\mu^\star$ returns the same first-stage decision
as the full smoothed stochastic program.
\end{theorem}

\begin{proof}
\emph{Step 1: $\bar\lambda_\mu$ is strictly dual-feasible for $q^\star$.}
For each scenario $i$, Proposition~\ref{prop:second-stage-kkt} gives
\[
q_i-W^\top\lambda_{\mu}^{\star,i}=\mu\,\bigl(y_{\mu,i}^\star\bigr)^{-1}\in\mathbb{R}^{n_2}_{++}.
\]
By \eqref{eq:qmax}, $q^\star\ge q^{\max}\ge q_i$ componentwise, so
$q^\star-W^\top\lambda_{\mu}^{\star,i}\in\mathbb{R}^{n_2}_{++}$ for all $i$.
Taking the $p_i$-weighted average yields
\[
q^\star-W^\top\bar\lambda_\mu
=\sum_{i=1}^N p_i\bigl(q^\star-W^\top\lambda_{\mu}^{\star,i}\bigr)\in\mathbb{R}^{n_2}_{++}.
\]

\emph{Step 2: Construct $(y_\mu^\star,h_\mu^\star)$ so that $\bar\lambda_\mu$ is second-stage optimal.}
Define $y_\mu^\star:=\mu(q^\star-W^\top\bar\lambda_\mu)^{-1}$ and $h_\mu^\star:=Tz_\mu^\star+Wy_\mu^\star$.
Then $Wy_\mu^\star=h_\mu^\star-Tz_\mu^\star$ by construction, and also
$q^\star-W^\top\bar\lambda_\mu=\mu(y_\mu^\star)^{-1}$.
Thus the pair $(y_\mu^\star,\bar\lambda_\mu)$ satisfies the KKT system \eqref{eq:KKT-second-stage}
for the \emph{single-scenario} second-stage barrier problem with data $(q^\star,W,T,h_\mu^\star)$ at $z=z_\mu^\star$.
In particular, $\bar\lambda_\mu\in\Lambda_\mu(z_\mu^\star;\hat\xi_\mu^\star)$.

\emph{Step 3: Verify first-stage optimality for the one-scenario problem.}
Because $T$ is fixed across scenarios, the first-stage KKT condition \eqref{eq:KKT-first-stage} for the
stochastic problem reads
\[
c-\mu(z_\mu^\star)^{-1}+A^\top\nu_\mu^\star - T^\top\Bigl(\sum_{i=1}^N p_i\lambda_{\mu}^{\star,i}\Bigr)=0,
\]
i.e.\ $c-\mu(z_\mu^\star)^{-1}+A^\top\nu_\mu^\star-T^\top\bar\lambda_\mu=0$.
Together with $Az_\mu^\star=b$ and the second-stage KKT relations from Step~2 (with multiplier
$\bar\lambda_\mu$), this is exactly the KKT system for the one-scenario barrier problem with scenario
$\hat\xi_\mu^\star$ at $z=z_\mu^\star$. Hence $z_\mu^\star$ is optimal for the one-scenario problem,
proving \eqref{eq:decision-optimal-qh}.
\end{proof}

\subsubsection[Noise in (W,q,h) with fixed T, and (W*, h*) chosen a priori]{Noise in $(W,q,h)$ with fixed $T$, and $(W^\star,h^\star)$ chosen \emph{a priori}}\label{subsec:Wqh-fixedT}

Assume now that $T_i\equiv T$ is fixed across scenarios, but $(W_i,q_i,h_i)$ may vary with $i$.

\begin{theorem}[A decision-optimal $q_\mu^\star$ under fixed $T$ and a priori $(W^\star,h^\star)$]\label{thm:Wqh-fixedT-barrier}
Assume $T$ is fixed across scenarios. Let $z_\mu^\star$ be the unique minimizer of the smoothed stochastic
program \eqref{eq:stoch-mu}, and define the averaged dual vector
\[
\bar\lambda_\mu:=\sum_{i=1}^N p_i\,\lambda_{\mu}^{\star,i}.
\]
Fix an index $i_0\in\{1,\dots,N\}$ (chosen a priori) and set
\[
W^\star:=W_{i_0},\qquad h^\star:=h_{i_0},\qquad \lambda_{0,\mu}:=\lambda_{\mu}^{\star,i_0}.
\]
Define
\begin{equation}\label{eq:qstar-Wqh}
q_\mu^\star:=q_{i_0}+{W^\star}^\top(\bar\lambda_\mu-\lambda_{0,\mu}),
\end{equation}
and the single-scenario forecast $\hat\xi_\mu^\star:=(q_\mu^\star,W^\star,T,h^\star)$.
Then
\[
z_\mu^\star = z_\mu^\star(\hat\xi_\mu^\star),
\]
i.e.\ $\hat\xi_\mu^\star$ is a decision-optimal single-scenario forecast for the barrier-smoothed problem.
\end{theorem}

\begin{proof}
\emph{Step 1: $\bar\lambda_\mu$ is second-stage optimal for $(q_\mu^\star,W^\star,h^\star)$ at $z_\mu^\star$.}
By definition, $\Delta:=\bar\lambda_\mu-\lambda_{0,\mu}$ and $q_\mu^\star=q_{i_0}+{W^\star}^\top\Delta$.
Lemma~\ref{lem:q-translation} applies with $W=W^\star$, $h=h^\star$ and $z=z_\mu^\star$, and yields
\[
\Lambda_\mu\bigl(z_\mu^\star;\,q_\mu^\star,W^\star,T,h^\star\bigr)
=
\Lambda_\mu\bigl(z_\mu^\star;\,q_{i_0},W^\star,T,h^\star\bigr)+\Delta.
\]
Since $\lambda_{0,\mu}\in\Lambda_\mu(z_\mu^\star;\,q_{i_0},W^\star,T,h^\star)$, it follows that
$\lambda_{0,\mu}+\Delta=\bar\lambda_\mu$ belongs to
$\Lambda_\mu(z_\mu^\star;\,q_\mu^\star,W^\star,T,h^\star)$.

\emph{Step 2: First-stage optimality.}
Since $T$ is fixed across scenarios, the smoothed first-stage KKT condition \eqref{eq:KKT-first-stage} can be written
\[
c-\mu(z_\mu^\star)^{-1}+A^\top\nu_\mu^\star - T^\top\bar\lambda_\mu=0.
\]
Together with feasibility $Az_\mu^\star=b$ and Step~1 (which provides a valid second-stage multiplier
$\bar\lambda_\mu$ for the one-scenario recourse at $z_\mu^\star$), this is precisely the KKT system of the
one-scenario barrier problem with scenario $\hat\xi_\mu^\star$. Hence $z_\mu^\star$ is optimal for the
one-scenario problem with $\hat\xi_\mu^\star$.
\end{proof}

\subsubsection[Full noise in (W,T,q,h), with (W*, T*) chosen a priori]{Full noise in $(W,T,q,h)$, with $(W^\star,T^\star)$ chosen \emph{a priori}}\label{subsec:full-noise}

In the fully stochastic case, the first-stage KKT condition involves the combination $\sum_i p_i T_i^\top\lambda_{\mu}^{\star,i}$,
which generally cannot be written as $T^\top(\cdot)$ for a \emph{fixed} $T$ unless $T$ is fixed.
We therefore choose an \emph{a priori} matrix $T^\star$ whose row space dominates all scenario row spaces.

\begin{assumption}[Row-space dominance for $T^\star$]\label{ass:rowspace}
There exists a matrix $T^\star\in\mathbb{R}^{m^\star\times n_1}$, chosen a priori, such that
\[
\Row(T_i)\subseteq \Row(T^\star)\quad\text{for all }i=1,\dots,N,
\]
and $T^\star$ has full row rank.
\end{assumption}

\begin{remark}[Constructing such a $T^\star$]\label{rem:construct-Tstar}
Because the support is finite, one can construct $T^\star$ by stacking all $T_i$ vertically, extracting a basis of the resulting
row space (e.g.\ by Gaussian elimination), and taking those basis vectors as the rows of $T^\star$.
\end{remark}

We also choose a compatible $W^\star$ by selecting the corresponding rows from the scenario matrices $W_i$.

\begin{assumption}[Compatible choice of $W^\star$]\label{ass:Wstar}
Let the rows of $T^\star$ be selected from the collection of scenario rows $\{(T_i)_{j,:}\}$.
For each row of $T^\star$, pick the matching row of $W$ from the same scenario, and assemble them into a matrix
$W^\star\in\mathbb{R}^{m^\star\times n_2}$. We assume (and enforce by removing redundant rows) that $W^\star$ has full row rank.
\end{assumption}

\begin{theorem}[A decision-optimal single scenario under full noise]\label{thm:full-noise-barrier}
Assume \ref{ass:rowspace}--\ref{ass:Wstar}. Let $z_\mu^\star$ be the unique minimizer of the smoothed stochastic program
\eqref{eq:stoch-mu}.
Define
\[
g_\mu:=\sum_{i=1}^N p_i\,T_i^\top\lambda_{\mu}^{\star,i}\in\mathbb{R}^{n_1}.
\]
Then $g_\mu\in\Row(T^\star)=\col({T^\star}^\top)$, so there exists a unique vector $\lambda_\mu^\star\in\mathbb{R}^{m^\star}$
satisfying
\begin{equation}\label{eq:lambda-star-full-noise}
{T^\star}^\top\lambda_\mu^\star=g_\mu.
\end{equation}
Fix any reference vector $y^{\mathrm{ref}}\in\mathbb{R}^{n_2}_{++}$ (chosen a priori) and define
\begin{equation}\label{eq:full-noise-forecast}
q_\mu^\star:={W^\star}^\top\lambda_\mu^\star+\mu\,(y^{\mathrm{ref}})^{-1},
\qquad
h_\mu^\star:=T^\star z_\mu^\star+W^\star y^{\mathrm{ref}}.
\end{equation}
Let $\hat\xi_\mu^\star:=(q_\mu^\star,W^\star,T^\star,h_\mu^\star)$. Then
\[
z_\mu^\star = z_\mu^\star(\hat\xi_\mu^\star),
\]
i.e.\ $\hat\xi_\mu^\star$ is a decision-optimal single-scenario forecast for the barrier-smoothed problem.
\end{theorem}

\begin{proof}
\emph{Step 1: Existence and uniqueness of $\lambda_\mu^\star$.}
Each vector $T_i^\top\lambda_{\mu}^{\star,i}$ lies in $\Row(T_i)$, hence $g_\mu$ is a convex combination of vectors in $\Row(T_i)$.
Assumption~\ref{ass:rowspace} gives $\Row(T_i)\subseteq\Row(T^\star)$ for all $i$, hence $g_\mu\in\Row(T^\star)=\col({T^\star}^\top)$.
Since $T^\star$ has full row rank, ${T^\star}^\top$ has full column rank, so \eqref{eq:lambda-star-full-noise} has a unique solution.

\emph{Step 2: $\lambda_\mu^\star$ is second-stage optimal for the constructed one-scenario data.}
By construction \eqref{eq:full-noise-forecast}, the pair $(y^{\mathrm{ref}},\lambda_\mu^\star)$ satisfies
\[
W^\star y^{\mathrm{ref}} = h_\mu^\star - T^\star z_\mu^\star,
\qquad
q_\mu^\star - {W^\star}^\top\lambda_\mu^\star = \mu\,(y^{\mathrm{ref}})^{-1}.
\]
Thus it satisfies the KKT system \eqref{eq:KKT-second-stage} for the one-scenario second-stage barrier problem with data
$(q_\mu^\star,W^\star,T^\star,h_\mu^\star)$ at $z=z_\mu^\star$. In particular,
$\lambda_\mu^\star\in\Lambda_\mu(z_\mu^\star;\hat\xi_\mu^\star)$.

\emph{Step 3: First-stage optimality.}
The first-stage KKT condition \eqref{eq:KKT-first-stage} for the stochastic problem gives a multiplier $\nu_\mu^\star$ such that
\[
c-\mu(z_\mu^\star)^{-1}+A^\top\nu_\mu^\star - \sum_{i=1}^N p_i T_i^\top\lambda_{\mu}^{\star,i}=0,
\]
i.e.\ $c-\mu(z_\mu^\star)^{-1}+A^\top\nu_\mu^\star-g_\mu=0$.
Using \eqref{eq:lambda-star-full-noise}, $g_\mu={T^\star}^\top\lambda_\mu^\star$, so we obtain
\[
c-\mu(z_\mu^\star)^{-1}+A^\top\nu_\mu^\star-{T^\star}^\top\lambda_\mu^\star=0,
\]
which is precisely the one-scenario KKT stationarity for $\hat\xi_\mu^\star$ at $z_\mu^\star$.
Therefore $z_\mu^\star$ is optimal for the one-scenario barrier problem with scenario $\hat\xi_\mu^\star$.
\end{proof}

\subsection{From barrier forecasts to LP forecasts}\label{subsec:existence:LP}

We now relate the $\mu$-decision-optimal forecasts of Section~\ref{subsec:existence:mu} to their unsmoothed
LP counterparts as $\mu\downarrow 0$. Recall from Definition~\ref{def:ac} that for a one-scenario LP
we use $z^\star(\xi)$ to denote the analytic-center selection of its (possibly non-unique) optimal first-stage decisions.

\paragraph{Notation: analytic-center dual multipliers for unsmoothed recourse.}
Fix a first-stage decision $z\in Z$ and a scenario $\xi_i=(q_i,h_i,W_i,T_i)$.
The dual of the unsmoothed recourse LP $Q(z,\xi_i)$ is
\[
\max_{\lambda\in\mathbb{R}^{m_2}}\ \{(h_i-T_i z)^\top\lambda:\ W_i^\top\lambda\le q_i\}.
\]
Let $\Lambda^{\mathrm{LP}}(z;\xi_i)$ denote its (possibly set-valued) optimal solution set.
Define the set of dual-slack coordinates that are not identically zero on the optimal face by
\[
\mathcal J^{\mathrm{dual}}(z;\xi_i)
:=\Bigl\{k:\ \exists \lambda\in\Lambda^{\mathrm{LP}}(z;\xi_i)\ \text{with}\ (q_i-W_i^\top\lambda)_k>0\Bigr\}.
\]
The \emph{analytic-center dual multiplier} is the unique vector $\lambda^{\star,i}(z)\in\Lambda^{\mathrm{LP}}(z;\xi_i)$ solving
\[
\lambda^{\star,i}(z)\in
\argmax_{\lambda\in\Lambda^{\mathrm{LP}}(z;\xi_i)}
\ \sum_{k\in\mathcal J^{\mathrm{dual}}(z;\xi_i)}\log\bigl((q_i-W_i^\top\lambda)_k\bigr),
\]
where the maximization is taken over the relative interior of $\Lambda^{\mathrm{LP}}(z;\xi_i)$ in its affine hull.

\begin{lemma}[Closed-graph lemma at $\mu=0$ for one-scenario forecasts]\label{lem:closed-graph}
Let $\mu_k\downarrow 0$ and let $\hat\xi_k\to\hat\xi$ be any convergent sequence of one-scenario data.
Assume that for each $k$, the one-scenario barrier problem \eqref{eq:barrier-one-scenario} with data $\hat\xi_k$
is feasible and has a unique minimizer, and denote its first-stage component by $z_{\mu_k}^\star(\hat\xi_k)$.
If
\[
z_{\mu_k}^\star(\hat\xi_k)\to \bar z,
\]
then $\bar z = z^\star(\hat\xi)$, i.e.\ the limit is the analytic-center first-stage decision of the
\emph{unsmoothed} one-scenario LP with data $\hat\xi$.
\end{lemma}

\begin{proof}
Fix any $k$ and consider the map $\xi\mapsto (z_{\mu_k}^\star(\xi),y_{\mu_k}^\star(\xi))$.
For $\mu_k>0$, Proposition~\ref{prop:IFT} implies this map is continuous on any set where feasibility holds, hence
\[
\bigl\|z_{\mu_k}^\star(\hat\xi_k)-z_{\mu_k}^\star(\hat\xi)\bigr\|\xrightarrow[k\to\infty]{}0.
\]
Therefore $z_{\mu_k}^\star(\hat\xi)=z_{\mu_k}^\star(\hat\xi_k)+o(1)\to\bar z$.
On the other hand, for fixed $\hat\xi$, Theorem~\ref{thm:central-path-limit} shows that
$z_{\mu_k}^\star(\hat\xi)\to z^\star(\hat\xi)$ as $\mu_k\downarrow 0$.
Hence $\bar z=z^\star(\hat\xi)$.
\end{proof}

By Lemma~\ref{lem:closed-graph}, if a family of one-scenario forecasts $\{\hat\xi_\mu^\star\}_{\mu>0}$ satisfies
$z_\mu^\star = z_\mu^\star(\hat\xi_\mu^\star)$ for all $\mu>0$ and along some sequence $\mu_k\downarrow 0$ we have
$\hat\xi_{\mu_k}^\star\to \hat\xi^\star$, then $\hat\xi^\star$ is LP
decision-optimal, i.e.\ $z^\star = z^\star(\hat\xi^\star)$.

\subsubsection[A concrete LP limit for the (q*, h mu*) construction]{A concrete LP limit for the $(q^\star,h_\mu^\star)$ construction}\label{subsec:qh-limit}

We now specialize to the fixed-$(W,T)$ setting of Theorem~\ref{thm:qh-fixedWT-barrier}.
In that setting, $\hat\xi_\mu^\star=(q^\star,W,T,h_\mu^\star)$ with fixed $q^\star$ and varying $h_\mu^\star$.

\begin{theorem}[Convergence of the $h_\mu^\star$ construction]\label{thm:hmu-converges}
Assume the fixed-$(W,T)$ regime of Theorem~\ref{thm:qh-fixedWT-barrier}, and let $q^\star$ be
defined by \eqref{eq:qmax} with some fixed $\delta>0$. Let $z_\mu^\star$ be the unique minimizer
of the smoothed stochastic program \eqref{eq:stoch-mu}, let $\bar\lambda_\mu$ be as in Theorem~\ref{thm:qh-fixedWT-barrier},
and let $h_\mu^\star$ be defined by that theorem. Assume $z_\mu^\star\to z^\star$ as $\mu\downarrow 0$.
Then
\[
h_\mu^\star \xrightarrow[\mu\downarrow 0]{} h^\star:=Tz^\star.
\]
In particular, any accumulation point of the forecasts $(q^\star,W,T,h_\mu^\star)$ is the LP decision-optimal forecast
$(q^\star,W,T,Tz^\star)$.
\end{theorem}

\begin{proof}
Write $y_\mu^\star=\mu(q^\star-W^\top\bar\lambda_\mu)^{-1}$ so that $h_\mu^\star=Tz_\mu^\star+Wy_\mu^\star$.
We first show $y_\mu^\star\to 0$.
For each scenario $i$, the second-stage KKT gives $W^\top\lambda_{\mu}^{\star,i}=q_i-\mu(y_{\mu,i}^\star)^{-1}\le q_i$ componentwise.
Taking the probability-weighted average yields
\[
W^\top\bar\lambda_\mu=\sum_{i=1}^N p_i W^\top\lambda_{\mu}^{\star,i} \le \sum_{i=1}^N p_i q_i \quad\text{componentwise}.
\]
Let $q^{\max}\in\mathbb{R}^{n_2}$ be the componentwise maximum of the $(q_i)_{i=1}^N$, so that
$q^{\max}_k=\max_i(q_i)_k$ and $q^\star=q^{\max}+\delta|q^{\max}|$ by \eqref{eq:qmax}.
Since $\sum_i p_i q_i \le q^{\max}$ componentwise, we obtain
\[
q^\star - \sum_{i=1}^N p_i q_i \ \ge\ q^\star-q^{\max} \ =\ \delta\,|q^{\max}| \ \ge\ 0
\qquad\text{(componentwise)}.
\]
Combining with $W^\top\bar\lambda_\mu \le \sum_i p_i q_i$ yields
\[
q^\star - W^\top\bar\lambda_\mu \ \ge\ \delta\,|q^{\max}|.
\]
Therefore, componentwise,
\[
0<y_\mu^\star=\mu(q^\star-W^\top\bar\lambda_\mu)^{-1}
\ \le\ \mu\,(\delta|q^{\max}|)^{-1}
\ \xrightarrow[\mu\downarrow 0]{}\ 0,
\]
and hence $Wy_\mu^\star\to 0$.
Therefore $Wy_\mu^\star\to 0$ and, using $z_\mu^\star\to z^\star$,
\[
h_\mu^\star=Tz_\mu^\star+Wy_\mu^\star \to Tz^\star+0 = Tz^\star.
\]
Since $\hat\xi_\mu^\star=(q^\star,W,T,h_\mu^\star)$ and we have shown $h_\mu^\star\to Tz^\star$,
we obtain $\hat\xi_\mu^\star\to \hat\xi^\star:=(q^\star,W,T,Tz^\star)$.
Because $z_\mu^\star=z_\mu^\star(\hat\xi_\mu^\star)$ for all $\mu>0$
(Theorem~\ref{thm:qh-fixedWT-barrier}), Lemma~\ref{lem:closed-graph} yields
$z^\star=z^\star(\hat\xi^\star)$, i.e.\ $\hat\xi^\star$ is LP decision-optimal.
\end{proof}

\subsubsection[A concrete LP limit for the (q mu*, W*, T, h*) construction]{A concrete LP limit for the $(q_\mu^\star,W^\star,T,h^\star)$ construction}\label{subsec:Wqh-limit-LP}

We now specialize to the fixed-$T$ setting of Theorem~\ref{thm:Wqh-fixedT-barrier}.
In that setting, $\hat\xi_\mu^\star=(q_\mu^\star,W^\star,T,h^\star)$ with fixed $(W^\star,h^\star)$ and varying $q_\mu^\star$.

\begin{theorem}[LP limit of the $q_\mu^\star$ construction under fixed $T$]\label{thm:Wqh-fixedT-LP}
Assume the fixed-$T$ regime of Theorem~\ref{thm:Wqh-fixedT-barrier}: $T_i\equiv T$ is fixed across scenarios,
while $(W_i,q_i,h_i)$ may vary. Fix an index $i_0\in\{1,\dots,N\}$ \emph{a priori} and set
\[
W^\star:=W_{i_0},\qquad h^\star:=h_{i_0}.
\]
Let $\hat\xi_\mu^\star=(q_\mu^\star,W^\star,T,h^\star)$ be the $\mu$-decision-optimal barrier forecast from
Theorem~\ref{thm:Wqh-fixedT-barrier}, i.e.\ with $q_\mu^\star$ defined by \eqref{eq:qstar-Wqh} and
$z_\mu^\star=z_\mu^\star(\hat\xi_\mu^\star)$ for all $\mu>0$.

Let $z^\star$ denote the canonical (analytic-center) optimal first-stage decision of the \emph{unsmoothed}
extensive-form LP. Under the standing assumptions, $z_\mu^\star\to z^\star$ as $\mu\downarrow 0$
(by applying Theorem~\ref{thm:central-path-limit} to the extensive-form LP).

For each scenario $i$, let $\lambda^{\star,i}:=\lambda^{\star,i}(z^\star)$ denote the analytic-center dual-optimal
solution of the \emph{unsmoothed} second-stage dual LP at first-stage decision $z^\star$ and scenario $\xi_i$
(as defined in the notation block at the start of Section~\ref{subsec:existence:LP}). Define the averaged dual vector
\[
\lambda^\star := \sum_{i=1}^N p_i\,\lambda^{\star,i}.
\]
Then the LP-limit second-stage cost vector is given explicitly by
\begin{equation}\label{eq:qstar-fixedT-LP}
q^\star
:= q_{i_0} - {W^\star}^\top\bigl(\lambda^{\star,i_0}-\lambda^\star\bigr)
\;=\; q_{i_0} + {W^\star}^\top\bigl(\lambda^\star-\lambda^{\star,i_0}\bigr),
\end{equation}
i.e.\ in the form requested,
\[
q^\star = q_i - W_i^\top(\lambda^{\star,i}-\lambda^\star)\quad\text{with }i=i_0.
\]
Moreover, $q_\mu^\star\to q^\star$ as $\mu\downarrow 0$, hence $\hat\xi_\mu^\star\to \hat\xi^\star:=(q^\star,W^\star,T,h^\star)$.
In particular, $\hat\xi^\star$ is decision-optimal for the \emph{unsmoothed} problem:
\[
z^\star = z^\star(\hat\xi^\star).
\]
\end{theorem}

\begin{proof}
By Theorem~\ref{thm:Wqh-fixedT-barrier}, we have $z_\mu^\star=z_\mu^\star(\hat\xi_\mu^\star)$ for all $\mu>0$.
By Theorem~\ref{thm:central-path-limit} applied to the extensive-form LP, $z_\mu^\star\to z^\star$ as $\mu\downarrow 0$.

Next, use the standard primal--dual central-path convergence for the \emph{second-stage} problems (equivalently, for the
barrier-of-the-dual formulation in Proposition~\ref{prop:barrier-dual}) together with the fact that the analytic-center dual
solution depends continuously on the right-hand side $h_i-Tz$ (in particular, on $z$ under fixed $T$): this yields
\[
\lambda_{\mu}^{\star,i}\longrightarrow \lambda^{\star,i}\qquad (i=1,\dots,N),
\]
and hence $\bar\lambda_\mu=\sum_i p_i\lambda_{\mu}^{\star,i}\to\lambda^\star$ and
$\lambda_{0,\mu}=\lambda_{\mu}^{\star,i_0}\to\lambda^{\star,i_0}$.
Taking $\mu\downarrow 0$ in the definition $q_\mu^\star=q_{i_0}+{W^\star}^\top(\bar\lambda_\mu-\lambda_{0,\mu})$
gives $q_\mu^\star\to q^\star$ with $q^\star$ as in \eqref{eq:qstar-fixedT-LP}.
Therefore $\hat\xi_\mu^\star\to\hat\xi^\star$.

Finally, apply the convergent-scenario lemma (Lemma~\ref{lem:closed-graph}) to the sequence
$\hat\xi_\mu^\star\to\hat\xi^\star$ and the corresponding identities $z_\mu^\star=z_\mu^\star(\hat\xi_\mu^\star)$.
We obtain $z^\star=z^\star(\hat\xi^\star)$, i.e.\ $\hat\xi^\star$ is LP decision-optimal.
\end{proof}

\subsubsection{A concrete LP limit for the full-noise construction}\label{subsec:full-noise-limit-LP}

We now specialize to the full-noise setting of Theorem~\ref{thm:full-noise-barrier}.
In that setting, $\hat\xi_\mu^\star=(q_\mu^\star,W^\star,T^\star,h_\mu^\star)$ with fixed $(W^\star,T^\star)$ and varying $(q_\mu^\star,h_\mu^\star)$.

\begin{theorem}[LP limit of the full-noise construction]\label{thm:full-noise-LP}
Assume the full-noise regime and the a priori choices $(W^\star,T^\star)$ of
Assumptions~\ref{ass:rowspace}--\ref{ass:Wstar}. Let $z_\mu^\star$ denote the unique minimizer of the smoothed stochastic
program \eqref{eq:stoch-mu}, and let $z^\star$ denote the canonical (analytic-center) optimal first-stage decision of the
\emph{unsmoothed} extensive-form LP. Under the standing assumptions, $z_\mu^\star\to z^\star$ as $\mu\downarrow 0$
(by applying Theorem~\ref{thm:central-path-limit} to the extensive-form LP).

Define the LP-limit right-hand side by
\begin{equation}\label{eq:hstar-full-noise-LP}
h^\star := T^\star z^\star.
\end{equation}

For each scenario $i$, let $\lambda^{\star,i}:=\lambda^{\star,i}(z^\star)$ denote the analytic-center dual-optimal solution
of the \emph{unsmoothed} second-stage dual LP at first-stage decision $z^\star$ and scenario $\xi_i$.
Define
\[
g^\star := \sum_{i=1}^N p_i\,T_i^\top \lambda^{\star,i}\in\mathbb{R}^{n_1}.
\]
By Assumption~\ref{ass:rowspace}, $g^\star\in\Row(T^\star)=\col({T^\star}^\top)$, hence there exists a unique
$\lambda^\star\in\mathbb{R}^{m^\star}$ satisfying
\begin{equation}\label{eq:lambdastar-full-noise-LP}
{T^\star}^\top \lambda^\star = g^\star,
\qquad\text{equivalently}\qquad
\lambda^\star = (T^\star)^{-T}\,g^\star,
\end{equation}
where $(T^\star)^{-T}$ denotes the (unique) left-inverse action implied by \eqref{eq:lambdastar-full-noise-LP}.

Let $\lambda^{\star,0}\in\mathbb{R}^{m^\star}$ denote the analytic-center dual-optimal solution of the \emph{reference}
second-stage dual LP defined by first-stage decision $z=z^\star$ and second-stage data
\[
(W,T,h,q)=(W^\star,T^\star,h^\star,0),
\]
i.e.\ the analytic-center selection from the optimal set of
\[
\max_{\lambda\in\mathbb{R}^{m^\star}}\ \{(h^\star-T^\star z^\star)^\top\lambda:\ {W^\star}^\top\lambda\le 0\}.
\]
Then the LP-limit second-stage cost vector is chosen explicitly as
\begin{equation}\label{eq:qstar-full-noise-LP}
q^\star := -{W^\star}^\top\bigl(\lambda^{\star,0}-\lambda^\star\bigr),
\end{equation}
and the corresponding one-scenario LP forecast is $\hat\xi^\star:=(q^\star,W^\star,T^\star,h^\star)$.

Moreover, the full-noise barrier-optimal forecasts from Section~\ref{subsec:full-noise} converge to $\hat\xi^\star$ as
$\mu\downarrow 0$, and $\hat\xi^\star$ is decision-optimal for the \emph{unsmoothed} problem:
\[
z^\star = z^\star(\hat\xi^\star).
\]
\end{theorem}

\begin{proof}
By the full-noise barrier theorem (Theorem~\ref{thm:full-noise-barrier}), for each $\mu>0$ there is a
$\mu$-decision-optimal one-scenario barrier forecast $\hat\xi_\mu^\star$ such that
$z_\mu^\star=z_\mu^\star(\hat\xi_\mu^\star)$.

By Theorem~\ref{thm:central-path-limit} applied to the extensive-form LP, $z_\mu^\star\to z^\star$ as $\mu\downarrow 0$,
hence (by definition) $h_\mu^\star:=T^\star z_\mu^\star \to h^\star=T^\star z^\star$.

Next, by standard primal--dual central-path convergence for the \emph{second-stage} problems, the canonical barrier dual
multipliers converge to analytic-center dual-optimal multipliers of the unsmoothed second-stage dual LPs; in particular,
for each scenario $i$ we have $\lambda_{\mu}^{\star,i}\to\lambda^{\star,i}$, hence
$g_\mu:=\sum_i p_iT_i^\top\lambda_{\mu}^{\star,i}\to g^\star$.
Since the linear map $\lambda\mapsto {T^\star}^\top\lambda$ is injective and $g_\mu,g^\star\in\col({T^\star}^\top)$,
the unique solutions $\lambda_\mu^\star$ of ${T^\star}^\top\lambda_\mu^\star=g_\mu$ converge to the unique solution
$\lambda^\star$ of \eqref{eq:lambdastar-full-noise-LP}.

Finally, use continuity of the \emph{analytic-center dual solution map} for the second-stage dual LP with respect to the
right-hand side $h$ (with fixed $W^\star,T^\star$ and fixed $q=0$): since $h_\mu^\star\to h^\star$, the corresponding
reference analytic-center dual optimizers satisfy $\lambda_{\mu}^{\star,0}\to\lambda^{\star,0}$. Taking limits in the
(full-noise) $q$-translation defining the one-scenario forecast yields $q_\mu^\star\to q^\star$ with $q^\star$ given by
\eqref{eq:qstar-full-noise-LP}. Therefore $\hat\xi_\mu^\star\to\hat\xi^\star$.

Apply Lemma~\ref{lem:closed-graph} to $\hat\xi_\mu^\star\to\hat\xi^\star$ and $z_\mu^\star=z_\mu^\star(\hat\xi_\mu^\star)$.
We obtain $z^\star=z^\star(\hat\xi^\star)$, i.e.\ $\hat\xi^\star$ is LP decision-optimal.
\end{proof}

\begin{lemma}[Continuity of the analytic-center solution map]\label{lem:ac-continuity}
Suppose $\xi_k\to\xi$ and that the corresponding one-scenario LPs satisfy the standing assumptions guaranteeing
existence of analytic centers. Then $z^\star(\xi_k)\to z^\star(\xi)$.
\end{lemma}

\begin{proof}
Pick any sequence $\mu_\ell\downarrow 0$.
By Theorem~\ref{thm:central-path-limit}, for each fixed $k$ we have $z_{\mu_\ell}^\star(\xi_k)\to z^\star(\xi_k)$ as $\ell\to\infty$,
and similarly $z_{\mu_\ell}^\star(\xi)\to z^\star(\xi)$.
For each fixed $\mu_\ell>0$, continuity of $\xi\mapsto z_{\mu_\ell}^\star(\xi)$ (Proposition~\ref{prop:IFT}) yields
$z_{\mu_\ell}^\star(\xi_k)\to z_{\mu_\ell}^\star(\xi)$ as $k\to\infty$.
A standard diagonal argument then yields $z^\star(\xi_k)\to z^\star(\xi)$.
\end{proof}

\begin{remark}[From a barrier forecaster to a high-quality LP decision]\label{rem:procedure}
The results in this section justify the following practical pipeline for obtaining a high-quality \emph{unsmoothed} (LP) decision:
\begin{enumerate}
\item Choose a sufficiently small barrier parameter $\mu>0$.
\item Use the appropriate theorem from Section~\ref{subsec:existence:mu} to construct a $\mu$-decision-optimal
one-scenario barrier forecast $\hat\xi_\mu^\star$.
\item Solve the one-scenario \emph{barrier} problem with forecast $\hat\xi_\mu^\star$ to obtain its canonical barrier decision
$z_\mu^\star=z_\mu^\star(\hat\xi_\mu^\star)$.
\item Build a \emph{canonical LP forecast} from $z_\mu^\star$ using the corresponding theorem for the noise regime:
\begin{itemize}
\item fixed $(W,T)$: set $h:=Tz_\mu^\star$ (and keep the chosen $q^\star$);
\item fixed $T$ with a priori $(W^\star,h^\star)$: use the $q$-construction corresponding to Theorem~\ref{thm:Wqh-fixedT-LP};
\item full noise with a priori $(W^\star,T^\star)$: use the construction corresponding to Theorem~\ref{thm:full-noise-LP}.
\end{itemize}
Denote the resulting LP forecast by $\hat\xi^{\mathrm{LP}}_\mu$.
\item Solve the one-scenario \emph{unsmoothed} LP with forecast $\hat\xi^{\mathrm{LP}}_\mu$ and take the analytic-center decision
$z^\star(\hat\xi^{\mathrm{LP}}_\mu)$.
\end{enumerate}

If $\hat\xi^{\mathrm{LP}}_\mu\to \hat\xi^\star$ and $z_\mu^\star\to z^\star$ as $\mu\downarrow 0$, then
Lemma~\ref{lem:ac-continuity} implies $z^\star(\hat\xi^{\mathrm{LP}}_\mu)\to z^\star(\hat\xi^\star)=z^\star$.
\end{remark}
\bibliographystyle{plain}
% \bibliography{refs}

\end{document}
